% GENERATED FILE. Edit canonical lessons, then run npm run generate:books.
% canonical-source-hash: fnv1a64:b31398c3704c2e42
% canonical-lessons: KA-C170-herige, KA-C170-drsti, KA-C170-sonku, KA-C170-bokke, KA-C170-paustika

\chapter{Childbirth, Eyesight, Infection, Blister, Nutritious}
\label{ch:ka-childbirth-eyesight-infection-blister-nutritious}
\hlchaptermodality{170}

\begin{chapteropening}
\textbf{By the end of this chapter:} \emph{I can say childbirth, eyesight, an infection, a blister and nutritious.}

Complete the last lesson of chapter 170: I can say childbirth, eyesight, an infection, a blister and nutritious.
\end{chapteropening}

\section[\texorpdfstring{\kn{ಹೆರಿಗೆ} (herige)}{ಹೆರಿಗೆ (herige)}]{\kn{ಹೆರಿಗೆ} (herige) — childbirth, delivery}

\label{lesson:KA-C170-herige}



\begin{quote}
Before the new one: say the Kannada for curiosity, then the Kannada for a feeling.
\end{quote}

\subsection*{You'll want to know: \kn{ಹೆರಿಗೆ}}
\textbf{\kn{ಹೆರಿಗೆ}} — \emph{herige} — \textquotedblleft{}childbirth, delivery\textquotedblright{}.

\begin{grammarlens}[title={What you've built}]
One more word for this chapter.
\end{grammarlens}

\subsection*{Your turn}
\begin{itemize}
\raggedright
  \item \emph{Say it:} \emph{herige}
  \item \emph{Say it:} \emph{herige}, once more
  \item \emph{Say it:} say \emph{herige}
  \item \emph{Recall:} say the Kannada for curiosity, then the Kannada for a feeling, then say \emph{herige} again
\end{itemize}

\begin{tcolorbox}[breakable,colback=teal!4,colframe=teal!35!black,title={Before you move on}]
What does \kn{ಹೆರಿಗೆ} mean? (\textquotedblleft{}Childbirth, delivery\textquotedblright{}.) Say it once more.
\end{tcolorbox}
\section[\texorpdfstring{\kn{ದೃಷ್ಟಿ} (dṛṣṭi)}{ದೃಷ್ಟಿ (dṛṣṭi)}]{\kn{ದೃಷ್ಟಿ} (dṛṣṭi) — eyesight}

\label{lesson:KA-C170-drsti}



\begin{quote}
Before the new one: say the Kannada for a feeling, then the Kannada for childbirth.
\end{quote}

\subsection*{You'll want to know: \kn{ದೃಷ್ಟಿ}}
\textbf{\kn{ದೃಷ್ಟಿ}} — \emph{dṛṣṭi} — \textquotedblleft{}eyesight\textquotedblright{}.

\begin{grammarlens}[title={What you've built}]
One more word for this chapter.
\end{grammarlens}

\subsection*{Your turn}
\begin{itemize}
\raggedright
  \item \emph{Say it:} \emph{dṛṣṭi}
  \item \emph{Say it:} \emph{dṛṣṭi}, once more
  \item \emph{Say it:} say \emph{dṛṣṭi}
  \item \emph{Recall:} say the Kannada for a feeling, then the Kannada for childbirth, then say \emph{dṛṣṭi} again
\end{itemize}

\begin{tcolorbox}[breakable,colback=teal!4,colframe=teal!35!black,title={Before you move on}]
What does \kn{ದೃಷ್ಟಿ} mean? (\textquotedblleft{}Eyesight\textquotedblright{}.) Say it once more.
\end{tcolorbox}
\section[\texorpdfstring{\kn{ಸೋಂಕು} (sōṅku)}{ಸೋಂಕು (sōṅku)}]{\kn{ಸೋಂಕು} (sōṅku) — an infection}

\label{lesson:KA-C170-sonku}



\begin{quote}
Before the new one: say the Kannada for childbirth, then the Kannada for eyesight.
\end{quote}

\subsection*{You'll want to know: \kn{ಸೋಂಕು}}
\textbf{\kn{ಸೋಂಕು}} — \emph{sōṅku} — \textquotedblleft{}an infection\textquotedblright{}.

\begin{grammarlens}[title={What you've built}]
One more word for this chapter.
\end{grammarlens}

\subsection*{Your turn}
\begin{itemize}
\raggedright
  \item \emph{Say it:} \emph{sōṅku}
  \item \emph{Say it:} \emph{sōṅku}, once more
  \item \emph{Say it:} say \emph{sōṅku}
  \item \emph{Recall:} say the Kannada for childbirth, then the Kannada for eyesight, then say \emph{sōṅku} again
\end{itemize}

\begin{tcolorbox}[breakable,colback=teal!4,colframe=teal!35!black,title={Before you move on}]
What does \kn{ಸೋಂಕು} mean? (\textquotedblleft{}An infection\textquotedblright{}.) Say it once more.
\end{tcolorbox}
\section[\texorpdfstring{\kn{ಬೊಕ್ಕೆ} (bokke)}{ಬೊಕ್ಕೆ (bokke)}]{\kn{ಬೊಕ್ಕೆ} (bokke) — a blister}

\label{lesson:KA-C170-bokke}



\begin{quote}
Before the new one: say the Kannada for eyesight, then the Kannada for an infection.
\end{quote}

\subsection*{You'll want to know: \kn{ಬೊಕ್ಕೆ}}
\textbf{\kn{ಬೊಕ್ಕೆ}} — \emph{bokke} — \textquotedblleft{}a blister\textquotedblright{}.

\begin{grammarlens}[title={What you've built}]
One more word for this chapter.
\end{grammarlens}

\subsection*{Your turn}
\begin{itemize}
\raggedright
  \item \emph{Say it:} \emph{bokke}
  \item \emph{Say it:} \emph{bokke}, once more
  \item \emph{Say it:} say \emph{bokke}
  \item \emph{Recall:} say the Kannada for eyesight, then the Kannada for an infection, then say \emph{bokke} again
\end{itemize}

\begin{tcolorbox}[breakable,colback=teal!4,colframe=teal!35!black,title={Before you move on}]
What does \kn{ಬೊಕ್ಕೆ} mean? (\textquotedblleft{}A blister\textquotedblright{}.) Say it once more.
\end{tcolorbox}
\section[\texorpdfstring{\kn{ಪೌಷ್ಟಿಕ} (pauṣṭika)}{ಪೌಷ್ಟಿಕ (pauṣṭika)}]{\kn{ಪೌಷ್ಟಿಕ} (pauṣṭika) — nutritious}

\label{lesson:KA-C170-paustika}



\begin{quote}
Before the new one: say the Kannada for an infection, then the Kannada for a blister.
\end{quote}

\subsection*{You'll want to know: \kn{ಪೌಷ್ಟಿಕ}}
\textbf{\kn{ಪೌಷ್ಟಿಕ}} — \emph{pauṣṭika} — \textquotedblleft{}nutritious\textquotedblright{}.

\begin{grammarlens}[title={What you've built}]
One more word for this chapter.
\end{grammarlens}

\subsection*{Your turn}
\begin{itemize}
\raggedright
  \item \emph{Say it:} \emph{pauṣṭika}
  \item \emph{Say it:} \emph{pauṣṭika}, once more
  \item \emph{Say it:} say \emph{pauṣṭika}
  \item \emph{Recall:} say the Kannada for an infection, then the Kannada for a blister, then say \emph{pauṣṭika} again
\end{itemize}

\begin{tcolorbox}[breakable,colback=teal!4,colframe=teal!35!black,title={Before you move on}]
What does \kn{ಪೌಷ್ಟಿಕ} mean? (\textquotedblleft{}Nutritious\textquotedblright{}.) Say it once more.
\end{tcolorbox}
